\documentclass{article}

% if you need to pass options to natbib, use, e.g.:
%     \PassOptionsToPackage{numbers, compress}{natbib}
% before loading neurips_2025

% The authors should use one of these tracks.
% Before accepting by the NeurIPS conference, select one of the options below.
% 0. "default" for submission
 \usepackage{neurips_2025}
% the "default" option is equal to the "main" option, which is used for the Main Track with double-blind reviewing.
% 1. "main" option is used for the Main Track
%  \usepackage[main]{neurips_2025}
% 2. "position" option is used for the Position Paper Track
%  \usepackage[position]{neurips_2025}
% 3. "dandb" option is used for the Datasets & Benchmarks Track
 % \usepackage[dandb]{neurips_2025}
% 4. "creativeai" option is used for the Creative AI Track
%  \usepackage[creativeai]{neurips_2025}
% 5. "sglblindworkshop" option is used for the Workshop with single-blind reviewing
 % \usepackage[sglblindworkshop]{neurips_2025}
% 6. "dblblindworkshop" option is used for the Workshop with double-blind reviewing
%  \usepackage[dblblindworkshop]{neurips_2025}

% After being accepted, the authors should add "final" behind the track to compile a camera-ready version.
% 1. Main Track
 % \usepackage[main, final]{neurips_2025}
% 2. Position Paper Track
%  \usepackage[position, final]{neurips_2025}
% 3. Datasets & Benchmarks Track
 % \usepackage[dandb, final]{neurips_2025}
% 4. Creative AI Track
%  \usepackage[creativeai, final]{neurips_2025}
% 5. Workshop with single-blind reviewing
%  \usepackage[sglblindworkshop, final]{neurips_2025}
% 6. Workshop with double-blind reviewing
%  \usepackage[dblblindworkshop, final]{neurips_2025}
% Note. For the workshop paper template, both \title{} and \workshoptitle{} are required, with the former indicating the paper title shown in the title and the latter indicating the workshop title displayed in the footnote.
% For workshops (5., 6.), the authors should add the name of the workshop, "\workshoptitle" command is used to set the workshop title.
% \workshoptitle{WORKSHOP TITLE}

% "preprint" option is used for arXiv or other preprint submissions
 % \usepackage[preprint]{neurips_2025}

% to avoid loading the natbib package, add option nonatbib:
%    \usepackage[nonatbib]{neurips_2025}


\usepackage[utf8]{inputenc} % allow utf-8 input
\usepackage[T1]{fontenc}    % use 8-bit T1 fonts
%\usepackage{hyperref}       % hyperlinks
\usepackage{url}            % simple URL typesetting
\usepackage{booktabs}       % professional-quality tables
\usepackage{amsfonts}       % blackboard math symbols
\usepackage{nicefrac}       % compact symbols for 1/2, etc.
\usepackage{microtype}      % microtypography
\usepackage{xcolor}         % colors


%---------------------------------------------------------------------------%
\usepackage{graphicx}%,subfigure}
\usepackage{wrapfig}
\usepackage{booktabs}
\usepackage{natbib}
\usepackage[colorlinks = true, citecolor = black, urlcolor  = blue]{hyperref}
%our imports

\usepackage{amsmath}
\usepackage{cleveref}
\usepackage{multirow}
\usepackage{subcaption}
\captionsetup[table]{skip=6pt}


%--------------------------------- new commands ---------------------------%

%norm symbol
\newcommand{\norm}[1]{\left \lVert #1 \right \rVert}

%nested itemization
\renewcommand{\labelitemii}{$\star$}
% comments of the authors (create your own)

\newcommand{\diego}[1]{\textcolor{orange!90!red}{//DD: #1//}}

%-----------------------------------------------------------%




\title{Title}

\newcommand{\samethanks}[1][\value{footnote}]{\footnotemark[#1]}

% 
%%%%%%%%% AUTHORS - PLEASE UPDATE
\author{
}


\begin{document}
\maketitle

\begin{abstract}
\end{abstract}

\section{MTEB (Eng. v2)}
\subsection{References}

\begin{table}[ht]
\centering
\resizebox{\columnwidth}{!}{
\begin{tabular}{lccc@{\hskip\tabcolsep\vrule width 0.1pt\hskip\tabcolsep}ccccccc}
\toprule
Model & Size & Avg (Task) & Avg (Type) & Retrieval & Rerank. & STS & Summ. & PairClass. & Class. & Clust. \\
\midrule        
EmbeddGemma measured (me) & 0.3B & 68.1 & 62.2  & 51.94  & 36.70 & 83.15 & 37.58 & 84.92 & 87.74 &  53.40 \\
EmbeddGemma measured (mteb) & 0.3B & 68.1 & 62.2  & 51.94  & 36.70 & 83.15 & 37.58 & 84.92 & 87.74 &  53.40 \\
EmbeddGemma reported & 0.3B & 69.7 & 65.1  & 55.7 & 47.4 & 83.6 & 37.6 & 87.3 & 87.6  & 56.6 \\


\midrule
Qwen3-emebd measured (me)  & 0.6B   &  &  &  &  &  &  &  &  &  \\
Qwen3-emebd measured (mteb)  & 0.6B   &  &  &  &  &  &  &  &  &  \\
Qwen3-emebd (reported)  & 0.6B   &  &  &  &  &  &  &  &  &  \\

Qwen3-emebd measured (me)  & 8B   & 73.60 & 66.37  & 67.74  & 44.87 & 82.94 & 33.94 & 87.94 & 90.34 & 56.83 \\
Qwen3-emebd measured (mteb)  & 8B   &  &  &  &  &  &  &  &  &  \\
Qwen3-emebd (reported)  & 8B   & 75.23 & 68.71 & 69.44 & 51.56 & 88.58 & 34.83 & 87.52 & 90.43  & 58.57 \\

\midrule
\multirow{2}{*}{F2LLMv1}       
    & 1 &  &  &  &  &  &  &  &  &  \\
    & 2 &  &  &  &  &  &  &  &  &  \\

\bottomrule 
\end{tabular}
}
\caption{\textbf{F2LLM v1 training data}.}
\end{table}





\label{intro}

\begin{table}[ht]
\centering
\resizebox{\columnwidth}{!}{
\begin{tabular}{lccc@{\hskip\tabcolsep\vrule width 0.1pt\hskip\tabcolsep}ccccccc}
\toprule
Model & Train run & Mean(Task) & Mean(Type) & Retrieval & Reranking & STS & Summariz. & PairClass. & Classification & Clustering \\
\midrule        
\multirow{2}{*}{t5gemma-2 encoder} 
    & 1 &  &  &  &  &  &  &  &  &  \\
    & 2 &  &  &  &  &  &  &  &  &  \\                                
\midrule
\multirow{2}{*}{t5gemma-2-decoder}  
    & 1 &  &  &  &  &  &  &  &  &  \\
    & 2 &  &  &  &  &  &  &  &  &  \\
\midrule
\multirow{2}{*}{gemma3}       
    & 1 &  &  &  &  &  &  &  &  &  \\
    & 2 &  &  &  &  &  &  &  &  &  \\
\midrule
\multirow{2}{*}{qwen3}    
    & 1 &  &  &  &  &  &  &  &  &  \\
    & 2 &  &  &  &  &  &  &  &  &  \\
\bottomrule 
\end{tabular}
}
\caption{\textbf{F2LLM v1 training data}.}
\end{table}


\begin{table}[ht]
\centering
\resizebox{\columnwidth}{!}{
\begin{tabular}{lccc@{\hskip\tabcolsep\vrule width 0.1pt\hskip\tabcolsep}ccccccc}
\toprule
Model & Train run & Date/Hr & Loss & Lr & wd & ep & deepspeed & f2llm code & data curriculum &  \\
\midrule        
\multirow{2}{*}{t5gemma-2 encoder}   
    & 1 &  &  &  &  &  &  &  &  &  \\
    & 2 &  &  &  &  &  &  &  &  &  \\                                
\midrule
\multirow{2}{*}{t5gemma-2-decoder}  
    & 1 &  &  &  &  &  &  &  &  &  \\
    & 2 &  &  &  &  &  &  &  &  &  \\
\midrule
\multirow{2}{*}{gemma3}       
    & 1 &  &  &  &  &  &  &  &  &  \\
    & 2 &  &  &  &  &  &  &  &  &  \\
\midrule
\multirow{2}{*}{qwen3}    
    & 1 &  &  &  &  &  &  &  &  &  \\
    & 2 &  &  &  &  &  &  &  &  &  \\
\bottomrule 
\end{tabular}
}
\caption{\textbf{F2LLM v1 dataset: Hyperparameters}.}
\end{table}





\begin{table}[h!]
\centering
\begin{tabular}{llrr}
\hline
\textbf{Task Type} & \textbf{Stage} & \textbf{\# Tasks} & \textbf{\# Samples} \\
\hline
Retrieval & 1 (33\%) & 29 & 1,459,811 \\
Retrieval & 2 (67\%) & 29 & 2,963,903 \\
Reranking & 1 (33\%) & 1 & 6,549 \\
Reranking & 2 (67\%) & 1 & 13,298 \\
STS & 1 (33\%) & 3 & 1,829 \\
STS & 2 (67\%) & 3 & 3,715 \\
Summarization & 1 (33\%) & 2 & 93,846 \\
Summarization & 2 (67\%) & 2 & 190,537 \\
PairClassification & 1 (33\%) & 3 & 61,201 \\
PairClassification & 2 (67\%) & 3 & 124,260 \\
Classification & 1 (33\%) & 6 & 66,998 \\
Classification & 2 (67\%) & 6 & 136,033 \\
Clustering & 1 (33\%) & 18 & 267,955 \\
Clustering & 2 (67\%) & 18 & 544,053 \\
\hline
\end{tabular}
\caption{Task distribution across stages}
\label{tab:task_distribution}
\end{table}

\clearpage 

% =============================================================
% End-to-end full-data runs (submit_train_f2llm_torchrun, full data)
% =============================================================

\begin{table}[ht]
\centering
\small
\setlength{\tabcolsep}{4pt}
\caption{End-to-end full-data runs: hyperparameter setup. All runs share backbone \texttt{t5gemma-2-270m-270m}, $4{\times}$H100 GPUs, global batch $320$, max seq.\ length $1024$, $7$ hard negatives, $1$ epoch, lr $10^{-5}$, weight decay $0.01$.}
\label{tab:e2e-hparams}
\begin{tabular}{lll}
\toprule
ID & Configuration & Train loss (final) \\
\midrule
E1 & plain backbone, EOS pool (baseline) & in-batch 0.159 / hard 0.503 @ step 18500 \\
E2 & gated attn pool, $h{=}4$, CLS query & in-batch 0.183 / hard 0.554 @ step 18500 \\
E3 & attn pool, $h{=}4$, CLS query & in-batch 2.553 / hard 1.806 @ step 17000 \\
E4 & attn pool $h{=}1$, frozen backbone & in-batch 0.879 / hard 1.284 @ step 18500 \\
E5 & cross-attn pool, layers 3/8/13/18/23 & in-batch 0.541 / hard 1.211 @ step 5000 \\
E6 & LoRA CLS attn, $r{=}64$ & in-batch 0.447 / hard 1.030 @ step 8000 \\
\bottomrule
\end{tabular}
\end{table}

\begin{table}[ht]
\centering
\small
\setlength{\tabcolsep}{4pt}
\caption{End-to-end full-data runs: MTEB-eng-v2 task accuracies (\%) at end of epoch 1.}
\label{tab:e2e-perf}
\begin{tabular}{lrrrrrrrrr}
\toprule
ID & Retrieval & Reranking & STS & Summarization & PairClassification & Classification & Clustering & micro & macro \\
\midrule
E1 & 53.55 & 45.25 & 81.31 & 30.69 & 81.69 & 87.38 & 55.21 & 67.67 & 62.16 \\
E2 & 51.62 & 45.13 & 81.19 & 25.36 & 82.46 & 86.10 & 54.63 & 66.73 & 60.93 \\
E3 & 8.88 & 45.48 & 58.09 & 14.71 & 48.36 & 49.05 & 38.18 & 36.23 & 37.54 \\
E4 & 28.65 & 39.08 & 69.34 & 29.96 & 73.73 & 69.94 & 43.74 & 52.42 & 50.64 \\
E5 & 42.20 & 57.14 & 75.89 & 24.59 & 73.35 & 68.04 & 39.55 & 54.87 & 54.39 \\
E6 & 44.83 & 57.65 & 76.45 & 24.68 & 75.58 & 69.68 & 42.02 & 56.69 & 55.84 \\
\bottomrule
\end{tabular}
\end{table}


% =============================================================
% Two-stage full-data runs (submit_train_2stage_f2llm_torchrun)
% =============================================================

\begin{table}[ht]
\centering
\small
\setlength{\tabcolsep}{4pt}
\caption{Two-stage full-data runs: hyperparameter setup and stage-1 final loss. Stage 1 trains the pooling head with backbone frozen; stage 2 finetunes end-to-end from the stage-1 head. All runs: $4{\times}$H100, batch $320$, max seq.\ $1024$, $7$ hard negatives, $1$ epoch, S1 fraction $1.0$.}
\label{tab:s2-hparams}
\begin{tabular}{ll l l}
\toprule
ID & Configuration & S1 loss (final) & S2 loss (final) \\
\midrule
T1 & gated attn $h{=}4$: S1 mean (lr$10^{-4}$/wd$0$) $\to$ S2 CLS (lr$10^{-5}$/wd$0.01$) & in-batch 0.798 / hard 1.235 & in-batch 0.187 / hard 0.527 \\
T2 & gated attn $h{=}4$: S1 mean (lr$10^{-4}$/wd$0$) $\to$ S2 CLS (lr$2{\cdot}10^{-5}$/wd$0.01$) & in-batch 0.798 / hard 1.235 & in-batch 0.197 / hard 0.543 \\
T3 & gated attn $h{=}4$: S1 mean $\to$ S2 mean (lr$2{\cdot}10^{-5}$/wd$0.01$) & in-batch 0.798 / hard 1.235 & in-batch 4.908 / hard 1.972 \\
T4 & ungated attn $h{=}4$: S1 mean $\to$ S2 CLS (lr$2{\cdot}10^{-5}$/wd$0.01$) & in-batch 1.188 / hard 1.436 & in-batch 0.292 / hard 0.708 \\
T5 & gated attn $h{=}8$: S1 mean $\to$ S2 CLS (lr$2{\cdot}10^{-5}$/wd$0.01$) & in-batch 0.947 / hard 1.321 & in-batch 0.549 / hard 0.998 \\
\bottomrule
\end{tabular}
\end{table}

\begin{table}[ht]
\centering
\small
\setlength{\tabcolsep}{4pt}
\caption{Two-stage full-data runs: MTEB-eng-v2 task accuracies (\%) at end of stage 2 epoch 1.}
\label{tab:s2-perf}
\begin{tabular}{lrrrrrrrrr}
\toprule
ID & Retrieval & Reranking & STS & Summarization & PairClassification & Classification & Clustering & micro & macro \\
\midrule
T1 & 48.85 & 44.75 & 80.81 & 29.95 & 82.43 & 86.26 & 55.31 & 66.22 & 61.19 \\
T2 & 48.93 & 44.87 & 80.51 & 30.40 & 82.45 & 85.87 & 55.10 & 66.07 & 61.16 \\
T3 & 0.23 & 40.32 & 38.78 & 7.87 & 18.09 & 31.61 & 29.49 & 22.31 & 23.77 \\
T4 & 47.33 & 43.45 & 79.12 & 26.32 & 81.48 & 80.90 & 51.69 & 63.51 & 58.61 \\
T5 & 40.15 & 41.42 & 75.38 & 20.65 & 77.09 & 65.49 & 46.31 & 56.32 & 52.36 \\
\bottomrule
\end{tabular}
\end{table}


% =============================================================
% Task-specific runs (submit_train_f2llm_torchrun with TASK_TYPES,
% and submit_train_2stage_single_task_torchrun)
% =============================================================

\begin{table}[ht]
\centering
\small
\setlength{\tabcolsep}{4pt}
\caption{Task-specific runs: variants. All runs share batch $320$, max seq.\ $1024$, $1$ epoch, lr $2{\cdot}10^{-5}$, wd $0.01$. Stage-1 (for 2-stage) uses lr $10^{-4}$, wd $0$.}
\label{tab:ts-hparams}
\begin{tabular}{ll}
\toprule
ID & Configuration \\
\midrule
S1 & plain backbone (EOS), single-stage \\
S2 & gated attn pool $h{=}4$, mean, single-stage \\
S3 & gated attn pool $h{=}4$, CLS, single-stage \\
S4 & gated attn $h{=}4$, 2-stage: S1 mean $\to$ S2 CLS \\
\bottomrule
\end{tabular}
\end{table}

\begin{table}[ht]
\centering
\small
\setlength{\tabcolsep}{4pt}
\caption{Task-specific runs: per-task MTEB accuracy (\%) at end of epoch 1. Each cell is the score for a model finetuned only on that task's data.}
\label{tab:ts-perf}
\begin{tabular}{lrrrrrrr}
\toprule
ID & Retrieval & Reranking & STS & Summarization & PairClassification & Classification & Clustering \\
\midrule
S1 & 53.71 & 43.08 & 62.92 & 29.04 & 82.51 & 60.85 & 57.12 \\
S2 & -- & 36.96 & 34.61 & 10.66 & 77.56 & 32.32 & 36.61 \\
S3 & 50.86 & 38.25 & 25.84 & 5.89 & 56.89 & 33.29 & 54.24 \\
S4 & -- & 34.74 & 22.56 & 15.62 & 80.39 & 31.54 & 38.47 \\
\bottomrule
\end{tabular}
\end{table}

\begin{table}[ht]
\centering
\small
\setlength{\tabcolsep}{4pt}
\caption{Task-specific runs: final reported training loss (Avg/global, in-batch / hard) at the last logged step.}
\label{tab:ts-train-loss}
\begin{tabular}{llllllll}
\toprule
ID & Retrieval & Reranking & STS & Summarization & PairClassification & Classification & Clustering \\
\midrule
S1 & 0.135/0.391 & -- & -- & 0.163/0.180 & 0.408/0.666 & --/0.583 & --/0.698 \\
S2 & -- & -- & -- & 4.608/2.041 & 1.220/1.170 & --/0.811 & --/1.932 \\
S3 & 0.163/0.466 & -- & -- & 4.822/2.067 & 5.697/2.109 & --/0.802 & --/0.801 \\
S4 & -- & -- & -- & 3.132/1.543 & 0.988/1.011 & --/0.926 & --/1.878 \\
\bottomrule
\end{tabular}
\end{table}






\clearpage

% =============================================================
% Procrustes alignment of layer representations
% =============================================================

\section{Generalized Procrustes alignment before attention pooling}
\label{sec:procrustes}

We add a per-layer learnable orthogonal alignment between the encoder
and the attention-pooling head, following the Generalized Procrustes
Analysis (GPA) scheme of Sch\"onemann (1966) / Gower (1975) as adapted
by \citet{arxiv:2602.06205}. Each of the $K$ per-view representations
$H_k \in \mathbb{R}^{B\times D}$ that the pooling head normally
consumes is first rotated by its own orthogonal matrix
$\Omega_k \in O(D)$, applied right-multiplicatively as $H_k\Omega_k$.
The pooling head then attends over the rotated stack
$(H_1\Omega_1,\dots,H_K\Omega_K)$.

\paragraph{Parameterisation.}
The orthogonal constraint is enforced exactly via
\texttt{torch.nn.utils.parametrizations.orthogonal} (matrix-exponential
of a learned skew-symmetric matrix). The new
\texttt{ProcrustesAlignment} module in \texttt{models/modules.py}
holds $K$ independent linear layers with the parametrisation applied
on top. Initialisation is selectable: \texttt{identity} (the
parameter is initialised to $I$, so the alignment is a no-op at step
zero and finetuning starts from the un-aligned baseline) or
\texttt{random} (random orthogonal).

\paragraph{GPA self-supervised loss.}
The same module exposes
\[
  \mathcal{L}_{\text{GPA}}(H)
  = \frac{1}{B K D}\sum_{k=1}^{K}\norm{H_k\Omega_k - U}_F^2,
  \qquad
  U = \tfrac{1}{K}\sum_{k=1}^{K} H_k\Omega_k,
\]
i.e.\ the squared deviation of each rotated view from the per-sample
consensus $U$ implicit in standard GPA. Backpropagation through this
loss only touches the orthogonal matrices --- the encoder runs under
\texttt{torch.no\_grad()} inside \texttt{compute\_gpa\_loss}.

\paragraph{Integration into the three pooling heads.}
The alignment plugs in front of the pooling stage of all three T5Gemma2
attention-pooling variants. The per-head value of $K$ is
\[
  K =
  \begin{cases}
    L+1 & \text{(\texttt{HiddenPool}, \texttt{pooling\_mode} $\in$ \{mean, cls\})}\\
    2(L+1) & \text{(\texttt{HiddenPool}, \texttt{pooling\_mode} = both)}\\
    \lvert\text{cross\_attention\_layers}\rvert + 1 & \text{(\texttt{CrossAttention})}\\
    L & \text{(\texttt{LoRACLSAttn})}
  \end{cases}
\]
where $L$ is \texttt{num\_hidden\_layers}. Each forward was refactored
to expose an \texttt{\_encode\_views(\dots)} helper that returns the
$(B,K,D)$ stack consumed by the pool, with the optional
\texttt{procrustes} layer inserted immediately before pooling. A
\texttt{gpa\_loss} kwarg on the \texttt{forward} method routes calls
through the GPA loss path while keeping the DeepSpeed engine wrapper
in the loop (so backward sees the right parameter set).

\paragraph{Two-stage workflow.}
The motivation for the orthogonal alignment is to pre-train the
$\Omega_k$ on a frozen encoder before contrastive finetuning, so the
finetuning starts from already-aligned layer representations. The
training driver \texttt{train\_f2llm\_repro.py} exposes:

\begin{itemize}
  \item \texttt{--procrustes\_alignment} enable the alignment layer.
  \item \texttt{--procrustes\_init \{identity, random\}} controls
        $\Omega_k$ initialisation.
  \item \texttt{--procrustes\_pretrain} (implies
        \texttt{--procrustes\_alignment}) freezes every parameter
        except \texttt{model.lm.procrustes}, switches the loss to
        $\mathcal{L}_{\text{GPA}}$, runs a dedicated loop
        \texttt{pretrain\_procrustes(\dots)} that calls
        \texttt{model.lm(\dots, gpa\_loss=True)} per batch, and saves
        the pre-trained state to \texttt{output\_dir/model.pt}.
  \item \texttt{--procrustes\_lr} optional separate learning rate
        for the pre-training stage.
  \item \texttt{--procrustes\_checkpoint <path>} for the
        contrastive finetuning run: loads only the
        \texttt{procrustes.*} tensors from a prior pre-training
        \texttt{model.pt}, leaving every other weight at its default
        initialisation. The matrices then continue to be updated
        end-to-end together with the rest of the model.
\end{itemize}

\noindent The flag \texttt{--procrustes\_alignment} is also surfaced
in \texttt{utils/arguments.py} and threaded through \texttt{train.py}
so both training drivers can opt in. A validation check rejects
\texttt{--procrustes\_alignment} when none of \texttt{--attention\_pooling},
\texttt{--lora\_cls\_attn}, or \texttt{--cross\_attention} is set
(those are the heads that consume a per-view stack).

\paragraph{Files touched.}
\texttt{models/modules.py} (new \texttt{ProcrustesAlignment}),
\texttt{models/t5gemma2model.py} (three pooling heads refactored:
constructor flag, \texttt{\_encode\_views} helper, \texttt{gpa\_loss}
forward kwarg, \texttt{compute\_gpa\_loss} method; factory
\texttt{get\_model\_t5gemma2\_model} propagates the new kwargs),
\texttt{f2llm\_repro/model.py} (\texttt{F2LLMT5Gemma2} reads
\texttt{args.procrustes\_alignment} / \texttt{args.procrustes\_init}),
\texttt{train\_f2llm\_repro.py} (CLI flags, validation, freeze/load
branches, \texttt{pretrain\_procrustes} loop, save logic),
\texttt{utils/arguments.py} and \texttt{train.py} (flag pass-through).


% =============================================================
% Catalogue of report files still tracking the codebase
% =============================================================

\section{Report directory --- current contents}
\label{sec:reports-catalog}

The \texttt{reports\_bugs\_refactoring/} directory accumulates short
markdown / text reports documenting past bug fixes and refactors.
Two files (\texttt{tokenizers\_parallelism.md} and
\texttt{compile\_error\_solved.md}) were removed: the first
contradicts the current \texttt{train.py} setting
(\texttt{TOKENIZERS\_PARALLELISM="true"}), and the second largely
duplicated \texttt{inplace\_op\_breaks\_backward.md} with an
additional ``fused query+doc forward'' note that no longer matches
\texttt{train.py}. What remains is summarised below; each entry is
the gist of the report plus a pointer to the code path that still
matches it.

\paragraph{\texttt{bug\_arguana.md} --- corpus all-gather inside
\texttt{search()}.}
DDP retrieval scores were halving with $\text{world\_size}>1$ because
the corpus chunk encoded with \texttt{divided\_by\_chunks=True}
returned local-only embeddings; each rank scored against $1/W$ of the
corpus. The fix all-gathers \texttt{local\_corpus\_chunk} and
\texttt{local\_indices} inside the search loop (with
\texttt{[:chunk\_n]} trim to remove the padding the length-sorted
sampler introduces). Still present in
\texttt{inference/helpers.py:search()} around the
\texttt{dist.all\_gather} block.

\paragraph{\texttt{inplace\_op\_breaks\_backward.md} --- inplace
mutations on autograd tensors broke \texttt{torch.compile} backward.}
\texttt{pairwise\_dot\_squared} in both
\texttt{EmbeddingGemmaLossDistributed} and
\texttt{EmbeddingGemmaLossHardNegatives} mutated \texttt{dots\_sq} via
\texttt{fill\_diagonal\_} / \texttt{masked\_fill\_}; the fix replaces
these with \texttt{torch.eye(B, dtype=bool)} +
\texttt{masked\_fill(\dots, 0)} out-of-place. Still applied in
\texttt{utils/losses.py}. The \texttt{fill\_diagonal\_(False)} calls
on non-autograd boolean tensors are intentionally left in place
because they cannot trigger the compile failure.

\paragraph{\texttt{memory\_optimization\_changes.txt} --- two-phase
memory rework for BioASQ (14\,M rows).} The first phase introduced
\texttt{LazyCorpusDict} (\texttt{tasks/data\_helpers.py}) backing the
old \texttt{dict[str, dict[str, str]]} pattern with an \texttt{id}
$\to$ index lookup over the existing text/title arrays, plus
\texttt{del} + \texttt{gc.collect()} at seven hot spots in
\texttt{tasks/retrieval\_loaders.py}. The second phase replaced
\texttt{Dataset(pa.table(\dots))} with a write-IPC-then-mmap
pipeline in \texttt{tasks/load\_datasets.py} to avoid the
\texttt{combine\_chunks()} double-copy that pushed each process past
its 125\,GB budget. Estimated peak savings: 15--25\,GB.

\paragraph{\texttt{models\_report.txt} --- bug history of the Gemma3
pooling models and the seed of the T5Gemma2 pooling variants.}
Documents eight fixes in \texttt{models/gemma3textmodel.py} +
\texttt{models/gemma3model.py} (\texttt{mean\_pool} keepdim,
\texttt{GatedAttention} kwargs, missing value/output projections,
\texttt{output\_hidden\_states} not forced True,
\texttt{torch.cat} $\to$ \texttt{torch.stack}, function-name
collision\dots), plus the initial port to T5Gemma2
(\texttt{EmbeddingT5Gemma2}, \texttt{EmbeddingT5Gemma2HiddenPool},
\texttt{--cls\_query\_pooling} option). All of the bug fixes are
still in place; the CLS-query and HiddenPool architectures described
here are the same ones that now host the Procrustes alignment
introduced in Sec.~\ref{sec:procrustes}.

\paragraph{\texttt{mteb\_20task\_subset\_selection.md} --- the
``reduced'' MTEB-eng-v2 task list.} A hand-picked subset of 20 of the
40 MTEB-eng-v2 tasks (5 retrieval, 4 clustering, 1 reranking, 3 STS,
4 classification, 2 pair-classification, 1 summarisation) chosen to
cut combined embeddinggemma + qwen3 evaluation time by $\sim$75\,\%
($2181\text{s} \to 554\text{s}$) while keeping per-category averages
within $\pm 3\%$ of the full set. The selection is exposed as the
\texttt{mteb\_eng\_v2\_reduced} option in \texttt{tasks/} +
\texttt{get\_eval\_tasks(\dots)}.

\paragraph{\texttt{overhead\_materialization\_python.md} --- the
20-minute string materialisation bottleneck in qrels filtering and
hard-negative mining.} Calls like
\texttt{pd.Series(qrels\_dataset["query\_id"])} and
\texttt{set(corpus\_dataset["id"])} forced HuggingFace
\texttt{Dataset} to decode 14\,M-row Arrow columns into Python
strings, each materialisation costing $\sim$3--4 minutes (six call
sites $\to$ 20\,min total). The fix uses PyArrow compute
(\texttt{pc.is\_in}, \texttt{pc.all}, \texttt{pc.count\_distinct})
directly on \texttt{dataset.data.table}, plus
\texttt{column(\dots).to\_pylist()} caching in \texttt{mine\_one()}
to avoid the eight redundant materialisations across
\texttt{search()} + \texttt{get\_hard\_negatives()}. Expected total
runtime: 23\,min $\to$ 1--2\,min. Still applied in
\texttt{utils/create\_datasets.py} and
\texttt{hard\_negative\_mining.py}.

\paragraph{\texttt{report\_train.md} --- the hard-negatives training
pipeline.} Documents the parquet layout under
\texttt{datasets\_negatives/<model>/} (query/positive/up-to-24
negatives per row, organised by retrieval / NLI / STS subfolders),
the new CLI surface (\texttt{--negatives\_dir},
\texttt{--num\_hard\_negatives}, \texttt{--instruction\_template},
\texttt{--max\_query\_len}, \texttt{--max\_passage\_len}), the
\texttt{TrainTaskMetadata} / \texttt{FOLDER\_TO\_TASK} mapping that
picks the right instruction template per subfolder, and the
\texttt{collate\_fn\_with\_hard\_negatives} / training-loop changes
that forward $B + B + B\!\cdot\! N$ sequences per step into
\texttt{EmbeddingGemmaLossHardNegatives}. The pipeline is still the
backbone of \texttt{train.py}.

\paragraph{\texttt{sampler\_shuffling\_considerations.md} ---
\texttt{DatasetAwareSampler}.} A DDP-aware sampler in
\texttt{utils/dataloader\_helpers.py} that guarantees every batch
contains rows from a single dataset, with two strategies behind the
\texttt{--batch\_strategy} flag: \texttt{sequential} (one dataset
fully consumed before the next; only the dataset order is shuffled
per epoch) and \texttt{grouped} (round-robin batches across
datasets). Both pad each dataset to a multiple of
\texttt{batch\_size $\times$ num\_replicas} and interleave indices
across ranks so per-rank length distributions stay balanced. The
report also justifies why length-sorted order is kept inside each
dataset (padding-efficiency dominates over the shuffle regularisation
benefit in single-epoch contrastive training). The
\texttt{--batch\_strategy} options enumerated here
(\texttt{mixed}, \texttt{sequential}, \texttt{grouped}, plus the
later \texttt{f2llm\_multi}) are the ones currently accepted in
\texttt{utils/arguments.py}.

\paragraph{\texttt{split\_changes\_summary.txt} --- adding
\texttt{validation}/\texttt{dev} splits to training data.} Documents
16 dataset loaders that were widened from \texttt{train} to
\texttt{train+validation} (or \texttt{train+dev}), covering the NLI
loaders (snli, mnli, anli, xnli), STS (stsbenchmark), 10 retrieval
loaders (squad, cnndm, xsum, miracl, mrtydi, hotpotqa, fiqa2018,
fever, msmarco, nfcorpus), and \texttt{kalm\_tasks/all\_nli}. The
report also lists the datasets explicitly left untouched (no
\texttt{validation}/\texttt{dev} available, or an mteb-style split
without an evaluation \texttt{test} part). The widened splits are
still active in \texttt{tasks/}.

\paragraph{\texttt{train\_bugs.txt} --- training-script review
(2026-02-28).} Reports four bugs and several improvement
suggestions. The applied fixes are: Bug 1 ---
\texttt{batch = \{k: v.to(\dots) if isinstance(v, Tensor) else v\}}
to stop \texttt{collate\_fn\_with\_hard\_negatives}'s plain-int
\texttt{num\_hard\_negatives} from crashing in
\texttt{Trainer.train(\dots)}; Bug 4 ---
\texttt{np.arange(step, max\_train\_steps+1, step)} so the first
log/eval slot is reachable (\texttt{utils/helpers.py:get\_cpt\_steps}).
Bugs 2 (3-D doc-id mask that only masks the positive column) and
3 (skip the hard-negative forward when the distributed loss ignores
them) and the six suggestions (spherical-loss weighting, symmetric
loss, configurable similarity-margin, restricting the spherical
loss to positives, distributed loss + hard negatives, gradient
accumulation, shuffling) are tracked but not yet applied to the
current code; treat this file as a partial backlog.

\paragraph{\texttt{training\_improvements.md} --- GPU-utilisation
optimisations (270M, short sequences).} Removed the 30\,\% idle
fraction caused by per-op kernel-launch overhead via: (1) fused
\texttt{AdamW} (\texttt{utils/optimizer.py}, \texttt{fused=True}),
(2) merging the query and document forward passes into a single
\texttt{self.model(\dots)} call in \texttt{train.py}, (3)
\texttt{non\_blocking=True} on the H2D copies (combined with the
DataLoader's \texttt{pin\_memory=True}), (4) DDP
\texttt{static\_graph=True} +
\texttt{gradient\_as\_bucket\_view=True} so the bucket schedule is
not recomputed each step. Optimisations 1, 2, 3, 5 still hold;
optimisation 4 (\texttt{mode="max-autotune"} on
\texttt{torch.compile}) is documented but the current
\texttt{train.py} uses the default compile mode --- adopt
\texttt{max-autotune} or \texttt{reduce-overhead} explicitly to get
the CUDA-graph capture benefit described here.


\end{document}


